\section{Analysis of Provided Log Material}
\label{sec:log-analysis}

This is where the actual investigation begins. Everything from this point onward uses only the
exports listed in Section 1.4, not the sandbox from Sections 3--4, which has already served
its purpose.

\begin{warningbox}
A quick note on tool soundness before you start querying: \texttt{jq}, used as below, only reads
the input file and prints to your terminal, it never modifies the file in place. \\
This makes it compatible with forensic soundness \emph{if you follow the right procedure}: hash each JSON file
before you touch it (as required in Section 1.4), never redirect \texttt{jq}'s output back onto the original file, and work from a copy. Note, though, that this soundness comes from \emph{how you use the tool}, not from \texttt{jq} being a certified forensic tool in the way Autopsy or Volatility are; it is a generic, general-purpose JSON processor. Keep this distinction in mind for Section 8.
\end{warningbox}

\subsection{Initial Survey}

\begin{lstlisting}
jq '.Records[] | {eventTime, eventName, sourceIPAddress, userIdentity: .userIdentity.userName}' cloudtrail_events.json
\end{lstlisting}
This query extracts, from every event record, only the four fields most relevant for a first
pass: when it happened, what happened, from where, and by whom. Scan the output chronologically.

\begin{questionbox}
List, in order, every event associated with the \texttt{svc-billing-sync} role on 19 May 2026.
What is anomalous about the \texttt{sourceIPAddress} of the first event in that sequence, and
how would you have discovered this anomaly systematically rather than by simply reading the
note field provided here?
\end{questionbox}

\begin{answerbox}
\vspace{3cm}
\end{answerbox}

\subsection{The Missing Data-Plane Events}

Recall from lecture the distinction between \emph{management-plane} events (actions on the cloud resources themselves: creating a bucket, changing its policy, listing its contents) and \emph{data-plane} events (actions on the actual data inside those resources: reading or writing individual objects).\\ Whether data-plane events are logged at all is a separate, opt-in configuration choice on most cloud platforms, it does not come "for free" alongside
management-plane logging.

\begin{questionbox}
Between the two \texttt{PutBucketPolicy} events at \texttt{03:16:22Z} and \texttt{03:45:51Z}, no \texttt{GetObject} events appear anywhere in the log. Explain precisely \emph{why} this is possible in a real cloud environment, referencing the distinction between management-plane and data-plane events. Is the absence of these events itself evidence, and if so, of what?
\end{questionbox}

\begin{answerbox}
\vspace{3cm}
\end{answerbox}

\subsection{Cross-Validating the Finding with a Second, Independent Method}

Recall from lecture that, during acquisition, using more than one hashing algorithm at once strengthens confidence in an integrity claim: if two independent computations agree, an error or manipulation in just one of them becomes far less plausible.\\
The same logic applies to analysis, not just acquisition: a finding that depends entirely on a single tool's correct behaviour is weaker than one confirmed by two independent tools that parse the data in different ways.
\medskip

\noindent
Confirm the absence of \texttt{GetObject} events from Section 5.2 using a tool that does not depend on \texttt{jq}'s JSON parser at all:

\begin{lstlisting}
grep -c '"eventName": "GetObject"' cloudtrail_events.json
\end{lstlisting}
This performs a dumb, literal text search for the string, entirely independently of \texttt{jq}'s
JSON interpretation. The result should be \texttt{0}, agreeing with what you observed in Section
5.2.

\begin{questionbox}
The \texttt{grep} command above and the \texttt{jq} query in Section 5.1 use fundamentally
different methods to read the same file (naive text search vs. structured JSON parsing). Why
does their agreement here strengthen your finding more than simply re-running the exact same
\texttt{jq} query a second time would? Can you construct a (contrived) scenario in which
\texttt{jq} and \texttt{grep} would legitimately disagree on this kind of question, even with
perfectly valid, unaltered JSON?
\end{questionbox}

\begin{answerbox}
\vspace{3cm}
\end{answerbox}

\noindent
Since the log export itself has a blind spot exactly where we need it most, we turn to a second,
independent artefact: a billing/data-transfer report. Billing systems and audit-logging systems
are usually built and operated independently of one another inside a cloud provider, precisely
because they serve different purposes, which is exactly what makes cross-referencing them
useful here.

\begin{lstlisting}
jq '.hourly_data_transfer_out_MB[]' data_transfer_metrics.json
\end{lstlisting}

\begin{questionbox}
Using \texttt{data\_transfer\_metrics.json}, identify the anomalous hour and estimate, as
precisely as the data allows, the order of magnitude of data transferred out of the bucket
during the un-logged window. This source was \textbf{not} disabled by the same bucket-policy
manipulation that suppressed \texttt{GetObject} logging. Why not? What does this tell you
about designing a logging architecture that is resilient to a single point of compromise?
\end{questionbox}

\begin{answerbox}
\vspace{3cm}
\end{answerbox}

\subsection{Reconstructing the Attacker's Actions}

Before writing the timeline, look again at the two earliest events of 19 May, \texttt{AssumeRole} and \texttt{GetSessionToken}. Every event after them is attributed to the role \texttt{svc-billing-sync}, which is a legitimate service identity and would raise no suspicion on its own. These two events are where that role was taken up, and they carry a field the later ones do not foreground: \texttt{sourceIdentity}, recording which principal invoked the assumption.

\begin{lstlisting}
jq '.Records[] | select(.eventName=="AssumeRole" or .eventName=="GetSessionToken") | {eventTime, eventName, sourceIPAddress, arn: .userIdentity.arn, sourceIdentity: .userIdentity.sourceIdentity}' cloudtrail_events.json
\end{lstlisting}

\begin{questionbox}
Distinguish the two identities present in these records: the role under which the subsequent actions were performed, and the principal that assumed it.
\\
(a) Had \texttt{sourceIdentity} been absent from the export, what could you still have concluded about \emph{who} acted, and what would have become unprovable?
\\
(b) An assumed role is a legitimate mechanism, used here by the billing service every day. What, in this specific sequence, separates the 19 May assumption from ordinary use, and would that distinction survive if the attacker had operated from an IP address already associated with the service?
\end{questionbox}

\begin{answerbox}
\vspace{3.5cm}
\end{answerbox}

\begin{questionbox}
Write a short timeline (bullet list, with UTC timestamps) reconstructing the attacker's actions
from first anomalous login to policy restoration, explicitly marking which steps are supported
by direct log evidence and which are inferred from indirect/circumstantial evidence (Section
5.4). Why is this distinction important if this timeline is later presented in court?
\end{questionbox}

\begin{answerbox}
\vspace{3.5cm}
\end{answerbox}
